% ID: SEL-D0028-OA-E02
% ESTADO: cotejada. Segunda pasada visual completa y borrador compilado revisados.
% Fuente: Selectividad/matematicas_aplicadas/suplente junio/Examen-Matemáticas Aplicadas a las Ciencias Sociales.pdf
% SHA-256: a10146fce64a17e21450973c88d66d43571619781f76bc345b5a4b3dee67a16b
% Página: 1; opción A; ejercicio original 2.

La velocidad que lleva un móvil, en función del tiempo $t$, viene dada por la siguiente función:
\[
v(t)=\begin{cases}
7t^2 & \text{si }0\leq t<1\\[4pt]
2t+a & \text{si }1\leq t\leq5\\[4pt]
-t^2+12t+b & \text{si }5<t\leq10
\end{cases}
\]
\begin{enumerate}[label=\alph*)]
\item \textbf{(1 punto)} Determine $a$ y $b$ para que la función sea continua en los instantes $t=1$ y $t=5$.
\item \textbf{(1.5 puntos)} Para $a=5$ y $b=-20$, estudie la derivabilidad en los instantes $t=1$ y $t=5$. \textquestiondown{}En qué momento el móvil alcanza la velocidad máxima?
\end{enumerate}
